% GENERATED FILE. Edit canonical lessons, then run npm run generate:books.
% canonical-source-hash: fnv1a64:d9e6fd48439d7290
% canonical-lessons: HI-C514-thappa, HI-C514-retna, HI-C514-jarna, HI-C514-camkana, HI-C514-pighlana

\chapter{Metal and Gold}
\label{ch:hi-metal-and-gold}
\hlchaptermodality{514}

\begin{chapteropening}
\textbf{By the end of this chapter:} \emph{I can describe how metal is filed, set and polished.}

Complete the last lesson of chapter 514: I can describe how metal is filed, set and polished.
\end{chapteropening}

\section[\texorpdfstring{\hi{ठप्पा} (ṭhappā)}{ठप्पा (ṭhappā)}]{\hi{ठप्पा} (ṭhappā) — a stamp, a die}

\label{lesson:HI-C514-thappa}



\begin{quote}
Before the new one: say the Hindi for to thread (a needle), then the Hindi for to come unstitched.
\end{quote}

\subsection*{You'll want to know: \hi{ठप्पा}}
\textbf{\hi{ठप्पा}} — \emph{ṭhappā} — \textquotedblleft{}a stamp, a die\textquotedblright{}.

\begin{quote}
\hi{अंगूठी} \hi{पर} \hi{सुनार} \hi{का} \hi{ठप्पा} \hi{है।}
\end{quote}

\begin{quote}
\emph{Aṅgūṭhī par sunār kā ṭhappā hai.}
\end{quote}

\begin{quote}
The ring bears the goldsmith's stamp.
\end{quote}

Masculine.

\begin{cousinweb}
\textbf{\hi{ठप्पा}} (\emph{ṭhappā}) is a tadbhava, inherited from \emph{*ṭhapp}, \textquotedblleft{}strike, press\textquotedblright{}, a form CDIAL reconstructs from its descendants (CDIAL 5495). CDIAL gives the Hindi word as \textquotedblleft{}stamp, mould\textquotedblright{}.
\end{cousinweb}

\begin{grammarlens}[title={What you've built}]
One more word for describing what you lived and felt.
\end{grammarlens}

\subsection*{Your turn}
\begin{itemize}
\raggedright
  \item \emph{Say it:} \emph{ṭhappā}
  \item \emph{Say it:} \emph{ṭhappā}, once more
  \item \emph{Say it:} \emph{Aṅgūṭhī par sunār kā ṭhappā hai.} — The ring bears the goldsmith's stamp.
  \item \emph{Recall:} say the Hindi for to thread (a needle), then the Hindi for to come unstitched, then say \emph{ṭhappā} again
\end{itemize}

\begin{tcolorbox}[breakable,colback=teal!4,colframe=teal!35!black,title={Before you move on}]
What does \hi{ठप्पा} mean? (\textquotedblleft{}A stamp, a die\textquotedblright{}.) Say it once more.
\end{tcolorbox}
\section[\texorpdfstring{\hi{रेतना} (retnā)}{रेतना (retnā)}]{\hi{रेतना} (retnā) — to file (metal or wood)}

\label{lesson:HI-C514-retna}



\begin{quote}
Before the new one: say the Hindi for to come unstitched, then the Hindi for a stamp, a die.
\end{quote}

\subsection*{You'll want to know: \hi{रेतना}}
\textbf{\hi{रेतना}} — \emph{retnā} — \textquotedblleft{}to file (metal or wood)\textquotedblright{}.

\begin{quote}
\hi{लोहार} \hi{ने} \hi{चाबी} \hi{का} \hi{किनारा} \hi{रेता।}
\end{quote}

\begin{quote}
\emph{Lohār ne cābī kā kinārā retā.}
\end{quote}

\begin{quote}
The blacksmith filed the edge of the key.
\end{quote}

Past, with \hi{ने} (he or she …): \hi{उसने} \hi{कील} \hi{रेती}, \emph{usne kīl retī}.

\begin{cousinweb}
\textbf{\hi{रेतना}} (\emph{retnā}) is a tadbhava, inherited from Sanskrit \emph{rētra}, \textquotedblleft{}perfumed powder\textquotedblright{} (CDIAL 10816), the entry under which CDIAL lists Hindi \emph{ret}, \textquotedblleft{}sand, a file\textquotedblright{}, and \emph{retnā}, \textquotedblleft{}to file\textquotedblright{}.
\end{cousinweb}

\begin{grammarlens}[title={What you've built}]
One more word for describing what you lived and felt.
\end{grammarlens}

\subsection*{Your turn}
\begin{itemize}
\raggedright
  \item \emph{Say it:} \emph{retnā}
  \item \emph{Say it:} \emph{retnā}, once more
  \item \emph{Say it:} \emph{Lohār ne cābī kā kinārā retā.} — The blacksmith filed the edge of the key.
  \item \emph{Recall:} say the Hindi for to come unstitched, then the Hindi for a stamp, a die, then say \emph{retnā} again
\end{itemize}

\begin{tcolorbox}[breakable,colback=teal!4,colframe=teal!35!black,title={Before you move on}]
What does \hi{रेतना} mean? (\textquotedblleft{}To file (metal or wood)\textquotedblright{}.) Say it once more.
\end{tcolorbox}
\section[\texorpdfstring{\hi{जड़ना} (jaṛnā)}{जड़ना (jaṛnā)}]{\hi{जड़ना} (jaṛnā) — to set (a stone); to inlay}

\label{lesson:HI-C514-jarna}



\begin{quote}
Before the new one: say the Hindi for a stamp, a die, then the Hindi for to file (metal or wood).
\end{quote}

\subsection*{You'll want to know: \hi{जड़ना}}
\textbf{\hi{जड़ना}} — \emph{jaṛnā} — \textquotedblleft{}to set (a stone); to inlay\textquotedblright{}.

\begin{quote}
\hi{सुनार} \hi{ने} \hi{अंगूठी} \hi{में} \hi{एक} \hi{लाल} \hi{पत्थर} \hi{जड़ा।}
\end{quote}

\begin{quote}
\emph{Sunār ne aṅgūṭhī meṁ ek lāl patthar jaṛā.}
\end{quote}

\begin{quote}
The goldsmith set a red stone in the ring.
\end{quote}

Past, with \hi{ने} (he or she …): \hi{उसने} \hi{कंगन} \hi{में} \hi{पत्थर} \hi{जड़ा}, \emph{usne kaṅgan meṁ patthar jaṛā}.

\begin{cousinweb}
\textbf{\hi{जड़ना}} (\emph{jaṛnā}) is a tadbhava, inherited from \emph{*jaḍati}, \textquotedblleft{}joins\textquotedblright{}, a form CDIAL reconstructs from its descendants (CDIAL 5091).
\end{cousinweb}

\begin{grammarlens}[title={What you've built}]
One more word for describing what you lived and felt.
\end{grammarlens}

\subsection*{Your turn}
\begin{itemize}
\raggedright
  \item \emph{Say it:} \emph{jaṛnā}
  \item \emph{Say it:} \emph{jaṛnā}, once more
  \item \emph{Say it:} \emph{Sunār ne aṅgūṭhī meṁ ek lāl patthar jaṛā.} — The goldsmith set a red stone in the ring.
  \item \emph{Recall:} say the Hindi for a stamp, a die, then the Hindi for to file (metal or wood), then say \emph{jaṛnā} again
\end{itemize}

\begin{tcolorbox}[breakable,colback=teal!4,colframe=teal!35!black,title={Before you move on}]
What does \hi{जड़ना} mean? (\textquotedblleft{}To set (a stone); to inlay\textquotedblright{}.) Say it once more.
\end{tcolorbox}
\section[\texorpdfstring{\hi{चमकाना} (camkānā)}{चमकाना (camkānā)}]{\hi{चमकाना} (camkānā) — to polish, to make shine}

\label{lesson:HI-C514-camkana}



\begin{quote}
Before the new one: say the Hindi for to file (metal or wood), then the Hindi for to set (a stone).
\end{quote}

\subsection*{You'll want to know: \hi{चमकाना}}
\textbf{\hi{चमकाना}} — \emph{camkānā} — \textquotedblleft{}to polish, to make shine\textquotedblright{}.

\begin{quote}
\hi{मीरा} \hi{ने} \hi{पुरानी} \hi{चाँदी} \hi{की} \hi{थाली} \hi{चमकाई।}
\end{quote}

\begin{quote}
\emph{Mīrā ne purānī cā̃dī kī thālī camkāī.}
\end{quote}

\begin{quote}
Mira polished the old silver platter.
\end{quote}

Past, with \hi{ने} (he or she …): \hi{उसने} \hi{बर्तन} \hi{चमकाए}, \emph{usne bartan camkāe}.

\begin{cousinweb}
\textbf{\hi{चमकाना}} (\emph{camkānā}) is the causative of \emph{camaknā}, \textquotedblleft{}to glitter\textquotedblright{}, which CDIAL derives from \emph{*camakka}, a reconstructed form (CDIAL 4676): it means \textquotedblleft{}to make glitter\textquotedblright{}.
\end{cousinweb}

\begin{grammarlens}[title={What you've built}]
One more word for describing what you lived and felt.
\end{grammarlens}

\subsection*{Your turn}
\begin{itemize}
\raggedright
  \item \emph{Say it:} \emph{camkānā}
  \item \emph{Say it:} \emph{camkānā}, once more
  \item \emph{Say it:} \emph{Mīrā ne purānī cā̃dī kī thālī camkāī.} — Mira polished the old silver platter.
  \item \emph{Recall:} say the Hindi for to file (metal or wood), then the Hindi for to set (a stone), then say \emph{camkānā} again
\end{itemize}

\begin{tcolorbox}[breakable,colback=teal!4,colframe=teal!35!black,title={Before you move on}]
What does \hi{चमकाना} mean? (\textquotedblleft{}To polish, to make shine\textquotedblright{}.) Say it once more.
\end{tcolorbox}
\section[\texorpdfstring{\hi{पिघलाना} (pighlānā)}{पिघलाना (pighlānā)}]{\hi{पिघलाना} (pighlānā) — to melt (something)}

\label{lesson:HI-C514-pighlana}



\begin{quote}
Before the new one: say the Hindi for to set (a stone), then the Hindi for to polish, to make shine.
\end{quote}

\subsection*{You'll want to know: \hi{पिघलाना}}
\textbf{\hi{पिघलाना}} — \emph{pighlānā} — \textquotedblleft{}to melt (something)\textquotedblright{}.

\begin{quote}
\hi{सुनार} \hi{ने} \hi{आग} \hi{पर} \hi{चाँदी} \hi{पिघलाई।}
\end{quote}

\begin{quote}
\emph{Sunār ne āg par cā̃dī pighlāī.}
\end{quote}

\begin{quote}
The goldsmith melted the silver over the fire.
\end{quote}

Past, with \hi{ने} (he or she …): \hi{उसने} \hi{बर्फ़} \hi{पिघलाई}, \emph{usne barf pighlāī}.

\begin{cousinweb}
\textbf{\hi{पिघलाना}} (\emph{pighlānā}) is the causative of \emph{pighalnā}, \textquotedblleft{}to melt\textquotedblright{}, which CDIAL files under \emph{*praghirati}, a reconstructed form, noting the older causative \emph{paghilānā} (CDIAL 8486).
\end{cousinweb}

\begin{grammarlens}[title={What you've built}]
One more word for describing what you lived and felt.
\end{grammarlens}

\subsection*{Your turn}
\begin{itemize}
\raggedright
  \item \emph{Say it:} \emph{pighlānā}
  \item \emph{Say it:} \emph{pighlānā}, once more
  \item \emph{Say it:} \emph{Sunār ne āg par cā̃dī pighlāī.} — The goldsmith melted the silver over the fire.
  \item \emph{Recall:} say the Hindi for to set (a stone), then the Hindi for to polish, to make shine, then say \emph{pighlānā} again
\end{itemize}

\begin{tcolorbox}[breakable,colback=teal!4,colframe=teal!35!black,title={Before you move on}]
What does \hi{पिघलाना} mean? (\textquotedblleft{}To melt (something)\textquotedblright{}.) Say it once more.
\end{tcolorbox}
